% ============================================================
% fonts.tex — 字体设置（中文宋体，英文 Times New Roman）
% ============================================================

% --- 中文字体 ---
% 学校规范要求中文统一用宋体。标题、目录、封面标签等需要粗体处使用
% 宋体合成粗体，不再切换到黑体或楷体。
\IfFontExistsTF{Songti SC}{%
  \setCJKmainfont[AutoFakeBold=2.5, AutoFakeSlant=0.167]{Songti SC}
  \setCJKsansfont[AutoFakeBold=2.5, AutoFakeSlant=0.167]{Songti SC}
  \setCJKmonofont[AutoFakeBold=2.5, AutoFakeSlant=0.167]{Songti SC}
  \setCJKfamilyfont{zhsong}[AutoFakeBold=2.5, AutoFakeSlant=0.167]{Songti SC}
  \setCJKfamilyfont{zhhei}[AutoFakeBold=2.5, AutoFakeSlant=0.167]{Songti SC}
  \setCJKfamilyfont{zhkai}[AutoFakeBold=2.5, AutoFakeSlant=0.167]{Songti SC}
}{%
  \setCJKmainfont[AutoFakeBold=2.5, AutoFakeSlant=0.167]{FandolSong}
  \setCJKsansfont[AutoFakeBold=2.5, AutoFakeSlant=0.167]{FandolSong}
  \setCJKmonofont[AutoFakeBold=2.5, AutoFakeSlant=0.167]{FandolSong}
  \setCJKfamilyfont{zhsong}[AutoFakeBold=2.5, AutoFakeSlant=0.167]{FandolSong}
  \setCJKfamilyfont{zhhei}[AutoFakeBold=2.5, AutoFakeSlant=0.167]{FandolSong}
  \setCJKfamilyfont{zhkai}[AutoFakeBold=2.5, AutoFakeSlant=0.167]{FandolSong}
}

\providecommand{\songti}{\CJKfamily{zhsong}}
\providecommand{\heiti}{\CJKfamily{zhhei}}
\providecommand{\kaishu}{\CJKfamily{zhkai}}

% --- 英文字体 ---
% 学校规范要求: 英文字体为 Times New Roman
\setmainfont{Times New Roman}
\setsansfont{Times New Roman}
\setmonofont{Times New Roman}
